% !TEX root = ../main.tex
\chapter*{Introduction générale}
\addcontentsline{toc}{chapter}{Introduction générale}
\markboth{Introduction générale}{}

% Plan conseillé : contexte général -> problème -> objectif du projet -> annonce du plan.
% Environ une page et demie. Pas de figure ni de tableau ici.

La transformation numérique du secteur de la santé s'accélère. Les patients attendent aujourd'hui de pouvoir réserver une consultation aussi simplement qu'un billet de train, à toute heure et sans passer par un appel téléphonique. Du côté des cliniques, la gestion des agendas reste pourtant souvent manuelle : registres papier, tableurs partagés et appels répétés au secrétariat.

Cette situation a un coût. Les secrétariats passent une grande partie de leur journée au téléphone, les praticiens subissent des créneaux vides lorsque les patients oublient leur rendez-vous, et la direction ne dispose d'aucun indicateur fiable sur l'occupation des cabinets.

C'est dans ce contexte que s'inscrit notre projet de fin d'études, réalisé au sein de \entreprise{} du \datedebut{} au \datefin. Il a pour objectif de concevoir et de réaliser \nomprojet, une plateforme SaaS qui permet à plusieurs cliniques de publier leurs disponibilités, aux patients de réserver en ligne, et à chaque établissement de suivre son activité grâce à un tableau de bord.

Le présent rapport s'articule autour de six chapitres :
\begin{itemize}
  \item le premier chapitre présente l'organisme d'accueil, le contexte du projet, l'étude de l'existant et la méthodologie adoptée ;
  \item le deuxième chapitre expose l'état de l'art, les solutions du marché et les choix technologiques ;
  \item le troisième chapitre détaille l'analyse et la spécification des besoins ;
  \item le quatrième chapitre décrit l'architecture et la conception de la solution ;
  \item le cinquième chapitre présente la réalisation, sprint par sprint ;
  \item le sixième chapitre traite des tests, de la qualité et du déploiement.
\end{itemize}
Nous clôturons ce rapport par une conclusion générale qui résume le travail réalisé et ouvre des perspectives.
